\section{Results}

\subsection{Does the stated intent decide the molecule?}
\label{sec:intent}

All arms answer the same \m{400} prompts from the molecule-level holdout, and the only thing that
differs between them is the sentence handed to $M$. Every comparison is paired on the prompt.

\begin{table}[H]\centering\small
\begin{tabular}{llrrrrr}
\toprule
arm & sentence given to $M$ & validity & uniq. & TC$\rightarrow$input & TC$\rightarrow$ref. (med.) & paired $\Delta$ vs mismatch \\
\midrule
reference (dictated) & agent-written, analogue in view & 0.963 & 0.626 & 0.677 & \textbf{0.918} (1.000) & $+0.379$, $p=5\!\times\!10^{-50}$ \\
$R\!\rightarrow\!M$ & $R$'s own generated intent & 0.965 & 0.238 & 0.746 & 0.570 (0.579) & $+0.033$, $p=0.014$ \\
mismatch & another pair's sentence & 0.915 & 0.937 & 0.718 & 0.537 (0.554) & --- (the control) \\
none & nothing & 1.000 & 0.235 & \textbf{1.000} & 0.684 (0.672) & $+0.147$, $p=5\!\times\!10^{-40}$ \\
\bottomrule
\end{tabular}
\caption{Intent conditioning on the molecule-level holdout, $n=400$ prompts per arm, PocketPair-21k.
Heavy-atom $\Delta$ MAE is \m{0.374} (oracle), \m{2.622} ($R\!\rightarrow\!M$) and \m{3.907}
(mismatch): under its own sentence the model moves about one heavy atom, under a foreign one it
drifts almost four. Validation-selection caveat applies.}
\label{tab:arms}
\end{table}

Three readings, in descending order of how much we would stake on them.

\textbf{A correct intent is worth $+0.379$ Tanimoto over a wrong one} (paired,
$p=5\times10^{-50}$), and the reference-intent median of \m{1.000} means that in more than half of
held-out cases the model returns the exact molecule that was made. This arm measures
\emph{transcription fidelity} and not design: the reference sentences were written by an agent with
the analogue in view, and they name a group that actually changed in \m{86.1\%} of rows (correctly
naming \m{2.09} of the \m{2.71} groups that change per pair). There is no oracle in this corpus. \textbf{A wrong intent is worse than
no intent}: the mismatched arm sits \m{0.147} below the no-intent arm ($p=5\times10^{-40}$) and its
validity falls from \m{100\%} to \m{91.5\%}. A model that merely decorated its output with a
rationale would be indifferent to which rationale it was handed; this one is not.

\textbf{The no-intent arm is a diagnostic, not a baseline.} With no sentence the model returns the
input unchanged in \m{100\%} of rows (TC$\rightarrow$input \m{1.000}), so its \m{0.684} is the
input's own similarity to the reference and all of its property deltas are exactly zero. This is an
artefact of pair training rather than a safety property --- trained on pairs with mostly small
edits, $p(y\mid x)$ with no instruction peaks at $y=x$, and the model learned \emph{copy, then
perturb where told}. It is nonetheless informative: zero intent, zero edit, which is direct evidence
that the sentence is the sole driver of structural change. The engineering consequence is that the
system has no sentence-free fallback mode.

\textbf{The reason model, not the molecule model, is the bottleneck.} $R\!\rightarrow\!M$ reaches
\m{0.570}: it clears the mismatched control only by \m{0.033} ($p=0.014$) and retains
$(0.570-0.537)/(0.918-0.537)=\m{8.7\%}$ of the oracle-over-mismatch gap. Its intents are also
repetitive (uniqueness \m{0.238} against oracle's \m{0.626}). In an earlier replication on the
crystal corpus the same ordering held with a smaller spread --- oracle \m{0.698}, $R\!\rightarrow\!M$
\m{0.518}, mismatch \m{0.503}, $+0.180$ at $p=3\times10^{-27}$, and invented-residue rates of
\m{0.0\%} (chemists), \m{13.1\%} ($R$), \m{27.1\%} ($R\!\rightarrow\!M$) and \m{55.7\%} (mismatch).
Two corpora, two splits, one ordering --- and two further points of arithmetic a reader will do
for us if we do not. First, $R\!\rightarrow\!M$ at \m{0.570} sits \emph{below} the no-intent arm at
\m{0.684}. That comparison is not meaningful, because the no-intent arm attains its figure by
returning the input unchanged, but it must be stated rather than left to be discovered. Second, the
two corpora report $R$ differently --- groundedness \m{0.869} and invented chemistry \m{13.1\%} with
residue specificity \m{0.421} on the crystal corpus, against groundedness \m{1.000} and specificity
\m{0.214} on the potency corpus. These are different corpora, different pockets and different
reference sentences, not two readings of one system: the crystal corpus has chemist-written
reference sentences to compare against, which is why its specificity is higher and its groundedness
lower. We report them side by side and never pool them. We therefore present the factorisation as a measured
\emph{upper bound} and as a product claim --- a chemist who can write the hypothesis gets the
structures today --- and not as an end-to-end gain. A direct comparison of the factorised pipeline
against a jointly trained model on one corpus and one split is \TBD{2}{joint vs.\ factorised};
hypothesis: the joint model matches $R\!\rightarrow\!M$ on molecular metrics and loses on every
rationale metric, because nothing in its objective makes the molecule depend on the reason.

\subsubsection{Beyond similarity: does a property-directed intent move that property?}

Tanimoto to a known answer cannot tell a chemist whether a stated goal was achieved. We therefore
classify each row by what its intent \emph{targets}, read from the sentence text rather than from the
assay, and compare the signed change (generated minus input) across arms on exactly those rows.

\begin{table}[H]\centering\small
\begin{tabular}{llrrrrl}
\toprule
the intent targets & property & rows & oracle & $R\!\rightarrow\!M$ & mismatch & verdict \\
\midrule
polarity $\uparrow$ & TPSA & 259 & $+3.62$ & $+4.84$ & $-0.12$ & moves as asked \\
lipophilicity $\downarrow$ & cLogP & 146 & $-0.63$ & $-0.36$ & $-0.16$ & moves as asked \\
size $\downarrow$ & MW & 75 & $-31.25$ & $+0.82$ & $-3.28$ & moves as asked (oracle only) \\
size $\uparrow$ & MW & 88 & $-1.98$ & $-2.66$ & $-10.54$ & flat \\
rigidity $\downarrow$ & rot.\ bonds & 37 & $+0.72$ & $-0.03$ & $+0.06$ & wrong direction \\
basicity $\downarrow$ & basic N & 31 & $-0.03$ & $-0.03$ & $+0.10$ & flat \\
\bottomrule
\end{tabular}
\caption{Property response, 349 paired rows. The no-intent arm is exactly $0.00$ in every cell
because it echoes the input, which is why it is omitted here. Bucket labels are read from the
conditioning sentence and are therefore partly circular with it; the informative comparison is
across arms within a row, not across rows.}
\label{tab:props}
\end{table}

Three of six buckets move in the stated direction and are flat or reversed under a mismatched
intent. Two fail outright. Steerability is real but uneven, and it is strongest exactly where a
single functional-group swap suffices --- adding polarity, shedding lipophilicity, trimming mass ---
and absent where the edit would have to be planned over several atoms. Confidence intervals and the
full bucket counts are \TBD{5}{property CIs}; hypothesis: the two failing buckets stay failing.

\subsection{What the pocket governs}
\label{sec:pocket}

The same corpus and the same held-out prompts, with one difference: whether the per-atom contact
table was present during training and generation.

\begin{table}[H]\centering\small
\begin{tabular}{lrrrr}
\toprule
 & groundedness & residue specificity & validity & TC$\rightarrow$reference \\
\midrule
contact table shown & \textbf{1.000} & \textbf{0.388} & 0.985 & 0.679 \\
contact table withheld & 0.432 & 0.121 & 1.000 & \textbf{0.680} \\
\bottomrule
\end{tabular}
\caption{Pocket ablation, $n=200$ held-out prompts per arm, PocketPair-21k.}
\label{tab:pocket}
\end{table}

The contact table decides \textbf{what the sentence may cite} --- groundedness \m{1.000} against
\m{0.432}, residue specificity \m{0.388} against \m{0.121} --- and it does \textbf{not} decide the
molecule: \m{0.679} with the table, \m{0.680} without it. Two further results point the same way. A
deliberately \emph{wrong} pocket, swapped in at cohort scale, costs \m{0.013} Tanimoto against the
true-analogue noise floor. And adding van der Waals distance annotations to the prompt changed
nothing measurable.

We report this as the central result rather than as a limitation, and we are explicit about what it
does and does not license. It does not license ``pocket-aware generation''. It does license
\emph{grounded attribution}: a target-specific, structurally checkable hypothesis attached to every
proposal, which is what a project team consumes when it has to defend a synthesis list.

Two caveats cut in opposite directions and both belong here. Against us: groundedness is close to
definitional, because the table sits in the prompt, so one arm can copy residue names and the other
must invent them; \m{1.000} versus \m{0.432} partly measures copying from context, and the only
non-trivial figure, residue specificity \m{0.388}, is low. A pocket-\emph{shuffled} control would
still score \m{1.000} by construction and is reported as such. In our favour: whole-molecule
Tanimoto is a poor instrument for this question. A single-edit analogue lives at
TC$\rightarrow$input \m{0.68}--\m{0.75}, where Tanimoto has almost no dynamic range, so
``the pocket did not move the molecule'' is partly a statement about the measurement. The
discriminating experiment is whether the edit lands on the atoms the table flags --- site-of-change
against cited-residue agreement --- and it is \TBD{12}{site-of-change vs.\ cited residue};
hypothesis: agreement exceeds a distance-sorted baseline but by less than the groundedness gap
suggests. Until that number exists, the negative result stands unmitigated.

\paragraph{A mechanism, and a fix.} The omission is specific and it is ours. Fields that describe
\emph{where a group could go} --- open room at an atom, slot shape, exit vectors --- are computed by
our structure layer and then deliberately excluded from the model's input, because they were
computed only at the edit site and therefore leaked it: the exit-vector field alone identified the
edit site in \m{93\%} of rows. The model is consequently handed a description of the current pose and
nothing about available space. A sentence can be grounded in that; a molecule cannot be derived from
it. Computing the same fields at \emph{every} heavy atom carries no site information and is
leakage-safe, and it is the one change we expect to make the pocket touch the molecule
\TBD{15}{room fields at every atom}; hypothesis: TC$\rightarrow$reference is unchanged but
site-of-change agreement improves.

\subsection{Generated molecules with a stated reason, and what finetuning buys}
\label{sec:gen}

\begin{table}[H]\centering\small
\begin{tabular}{llrrrrr}
\toprule
model & training data & validity & uniq. & TC$\rightarrow$input & sensible & sound reasoning \\
\midrule
ChemDFM-14B & none (base) & 0.118 & 0.871 & 0.793 & 10.6\% & 3.9\% \\
ether0-24B & none (base) & 0.225 & 0.847 & 0.884 & 12.7\% & 0.0\% \\
Llama-3.1-8B-Instruct & none (base) & 0.519 & 0.750 & 0.737 & 32.4\% & 6.5\% \\
SmileyLlama-8B & none (base) & 0.882 & 0.753 & 0.582 & 30.6\% & 0.0\% \\
\midrule
ChemDFM-14B & g2 (1{,}485) & 0.893 & 0.752 & 0.658 & 72.4\% & 22.9\% \\
ChemDFM-14B & g3 (2{,}326) & 0.958 & 0.777 & 0.683 & 95.8\% & 29.9\% \\
ChemDFM-14B & g4 (5{,}366) & 0.966 & 0.751 & 0.644 & 89.2\% & 37.6\% \\
ChemDFM-14B & PocketPair-9k (7{,}692) & \textbf{0.973} & 0.757 & 0.658 & 92.4\% & \textbf{46.5\%} \\
ChemDFM-14B & PocketPair-21k & 0.985 & \textbf{0.173} & 0.795 & \TBD{1}{} & \TBD{1}{} \\
\midrule
ether0-24B & PocketPair-9k & 0.855 & 0.768 & 0.639 & 79.6\% & 16.6\% \\
Llama-3.1-8B-Instruct & g3 & 0.771 & 0.743 & 0.594 & 57.8\% & 3.2\% \\
SmileyLlama-8B & g3 & 0.908 & 0.761 & 0.588 & 51.1\% & 3.8\% \\
DESRO-8B & PocketPair-14k & 0.954 & 0.756 & 0.662 & 90.0\% & 33.3\% \\
\bottomrule
\end{tabular}
\caption{Direct joint training. Panel columns come from one model-labeller panel under a fixed
rubric and seed; they are not comparable with panel figures elsewhere in this paper. Bases differ in
scale, so the only clean contrast is each base against itself.}
\label{tab:main}
\end{table}

Two trends and one cost. Validity \textbf{saturates} --- \m{0.893}, \m{0.958}, \m{0.966}, \m{0.973}
across a five-fold increase in data --- while sound reasoning \textbf{keeps climbing} ---
\m{22.9\%}, \m{29.9\%}, \m{37.6\%}, \m{46.5\%} --- and invented chemistry falls from \m{67.7\%} at
base to \m{42.0\%}. The chemistry is a plumbing problem solved early; the reasoning is the thing
data buys, and it had not plateaued when we stopped. The cost is diversity: the potency-trained
model's uniqueness collapses to \m{0.173}, about seventeen distinct ideas per hundred requests, and
$R$'s intents sit at \m{0.238}. On that evidence the system is a renderer of hypotheses, not a source
of them, and we say so rather than reporting uniqueness as a success.

Three numbers in this table will not carry a claim and we flag them here. \m{46.5\%} sound reasoning
and \m{42.0\%} invented chemistry are not ``usable'' by any standard a chemist would accept. The
panel columns rest on model labellers whose drift exceeds these differences. And every row is a
single run; seeded replicates with intervals are \TBD{16}{seed replication}.

\subsection{Agent-written versus literature-extracted reasoning}
\label{sec:agentic}

The intents supervising the crystal and potency corpora were written by an agent with the analogue
in view; the literature intents were written by the authors of the paper. We hold three results
about the former and none yet about the latter. On a cohort round of \m{613} candidates over five
targets, model labellers rated \m{0.61} on ``would make'' but only \m{0.08} on ``explanation
sound'': the molecules were acceptable and the reasons were not. An ownership audit scored \m{2/6}
where a panel scored \m{6/6}. And the field labelled ``specific'' in the agent's intent schema
turned out to measure the template rather than the content ($p=1.000$ on usefulness). Taken
together these say that agent-written reasoning is a weak teacher, which is the motivation for the
literature corpus rather than an aside.

What the literature corpus delivers is \TBD{1}{chemist panel on $R$'s intents} and
\TBD{2}{literature-trained $M$/$R$/$R\!\rightarrow\!M$}; hypotheses: $R$ trained on extracted
intents scores at parity with the agent-trained $R$ on groundedness and above it on usefulness,
while molecular metrics move little, because groundedness is already saturated by the contact table
and usefulness is the dimension the agent labels never had.

\subsection{Retrospective hit-to-lead}
\label{sec:retro}

With a publication-date cutoff and held-out series, the system proposes and justifies analogues of a
known binder and is scored on recovery of the compounds subsequently reported, on agreement with the
authors' own stated intent, and on potency improvement among analogues that were in fact measured,
against both a pocket-blind and a plain enumeration baseline. This experiment is
\TBD{3}{publication-date retrospective} and is gated on the complexes now being folded; it is also
our contamination control, since the base models were pretrained on literature and on ChEMBL that
plausibly include these pairs. Hypothesis: $R\!\rightarrow\!M$ beats the pocket-blind baseline on
intent agreement and lands between enumeration and oracle on recovery.

One caveat carries over from an earlier split study and must be read before any recovery number:
$R$ is the bottleneck by $\m{18.6\times}$, splitting costs roughly $\m{2\times}$ regardless, and the
figures of \m{0.188}/\m{0.703} from that study were \emph{recovery-near}, not exact recovery.

\subsection{A worked covalent programme: ZAP70}
\label{sec:zap70}

On one real covalent programme --- ZAP70, Cys346, \m{280} measured molecules --- a predictor over the
same structural features reaches absolute MAE \m{0.489} and Spearman \m{0.792} with docking into
4K2R plus pretraining, where docking and pretraining contribute additively, against an honestly
rebuilt baseline of MAE \m{0.723}. Earlier versions of the same pipeline reached \m{0.555} and
\m{0.554} (Spearman \m{0.726}). We note two things about this section rather than burying them. It
is a regression model, not $R$/$M$, and therefore evidence about the structural features rather than
about the factorisation. And a previously circulated figure of \m{0.469} for this target came from
one-leg pair leakage worth \m{0.163} MAE: it is withdrawn, it appears nowhere in this paper, and the
honest number is \m{0.723} for the baseline and \m{0.489} for the best model. At \m{280} molecules
every more complex architecture we tried made things worse, which is the most directly actionable
finding we have for programme-scale data: add structure, not parameters.
